\newpage
\section{INTRODUÇÃO}
O envelhecimento da população tem ampliado a necessidade de soluções que apoiem a manutenção da autonomia, do bem-estar e da integração social das pessoas idosas. 
Muitas organizações comunitárias e ONGs promovem atividades físicas, cognitivas e sociais voltadas a esse público, porém enfrentam desafios para manter a participação contínua dos membros, especialmente quando as interações dependem de encontros presenciais. Nesse contexto, surge a oportunidade de desenvolver uma plataforma digital que permita que pessoas idosas participem de atividades online e interajam com outros membros da comunidade.
A solução deve considerar características importantes desse público, como facilidade de uso, clareza visual, entre outras, promovendo experiências positivas de interação, reduzindo barreiras tecnológicas e fortalecendo o senso de pertencimento à comunidade. 
Além dos participantes idosos, há também o papel de mediadores comunitários (tutores), responsáveis por organizar atividades, compartilhar conteúdos e estimular a participação dos membros. A tecnologia deve apoiar esses mediadores na gestão das atividades e na comunicação com o grupo, garantindo que as interações ocorram de forma organizada e
segura.
A partir dessa narrativa inicial, o seguinte trabalho explora possíveis funcionalidades, as necessidades dos usuários e os requisitos do sistema, considerando aspectos de usabilidade, acessibilidade, interação social, suporte tecnológico e gestão das atividades.